\documentclass[12pt]{article}
\usepackage{amsmath}
\usepackage{hyperref}
\usepackage{enumitem}

\title{End-to-End Fraud Detection Project}
\author{}
\date{}

\begin{document}
\maketitle

\section*{Step 1: Project Planning}
\begin{itemize}
    \item \textbf{Define Project Objectives:}
    \begin{itemize}
        \item Predict whether a transaction is fraudulent.
        \item Visualize transaction data on the website using D3.js.
        \item Host the fraud detection model as a service on the website.
    \end{itemize}
    \item \textbf{Gather Requirements:}
    \begin{itemize}
        \item \textbf{Backend:} Python-based service for the fraud detection model.
        \item \textbf{Frontend:} HTML, CSS, JavaScript (with D3.js for visualization).
        \item \textbf{Deployment:} Cloud-hosted website (e.g., AWS, Azure, GCP, or Heroku).
    \end{itemize}
\end{itemize}

\section*{Step 2: Data Collection \& Preprocessing}
\begin{enumerate}
    \item Fetch data from GitHub stored in JSON format.
    \item Clean the data:
    \begin{itemize}
        \item Handle missing values.
        \item Remove or transform outliers.
        \item Standardize or normalize numerical features.
        \item Encode categorical variables.
    \end{itemize}
    \item Perform feature engineering:
    \begin{itemize}
        \item Create new features, e.g., transaction velocity, time since the last transaction.
        \item Remove irrelevant or redundant features.
    \end{itemize}
    \item Split the data into training, validation, and test sets.
\end{enumerate}

\section*{Step 3: Machine Learning Model}
\begin{enumerate}
    \item Select a binary classification algorithm such as Logistic Regression, Random Forests, or LightGBM.
    \item Train the model using the training data and perform hyperparameter tuning.
    \item Evaluate the model with metrics such as accuracy, precision, recall, F1-score, and ROC-AUC.
    \item Serialize the trained model using \texttt{pickle} or \texttt{joblib}.
\end{enumerate}

\section*{Step 4: Backend Development}
\begin{enumerate}
    \item Create a RESTful API for fraud prediction using \texttt{Flask} or \texttt{FastAPI}.
    \item Develop an endpoint (e.g., \texttt{/predict}) to accept transaction details and return predictions.
    \item Load the serialized model in the API and process user inputs.
    \item Test the API locally using tools like \texttt{Postman} or \texttt{cURL}.
\end{enumerate}

\section*{Step 5: Frontend Development}
\begin{enumerate}
    \item Design a simple HTML website with:
    \begin{itemize}
        \item A user input form for transaction details.
        \item Visualization of transaction data using D3.js.
        \item Display of fraud prediction results.
    \end{itemize}
    \item Use D3.js for data visualization:
    \begin{itemize}
        \item Load JSON data from GitHub using JavaScript.
        \item Create visualizations such as:
        \begin{itemize}
            \item Transaction timelines.
            \item Heatmaps of transaction locations.
            \item Bar charts for transaction frequency.
        \end{itemize}
    \end{itemize}
    \item Integrate the frontend with the backend using JavaScript (e.g., Fetch API or Axios).
\end{enumerate}

\section*{Step 6: Deployment}
\begin{enumerate}
    \item Deploy the backend API to a cloud platform (e.g., AWS Lambda, Heroku, or Google Cloud Run).
    \item Host the frontend (HTML/JavaScript) on a platform such as GitHub Pages, Netlify, or Vercel.
    \item Configure the frontend to communicate with the backend API.
\end{enumerate}

\section*{Step 7: Monitoring \& Scaling}
\begin{enumerate}
    \item Monitor model performance:
    \begin{itemize}
        \item Log predictions and user feedback.
        \item Periodically retrain the model with new data.
    \end{itemize}
    \item Scale the backend:
    \begin{itemize}
        \item Use load balancing or containerization (e.g., Docker and Kubernetes) for scaling.
    \end{itemize}
\end{enumerate}

\section*{Step 8: Documentation}
\begin{enumerate}
    \item Document the entire pipeline:
    \begin{itemize}
        \item Explain the data flow from GitHub to backend to frontend.
        \item Provide detailed documentation for API endpoints.
    \end{itemize}
    \item Write a \texttt{README.md} file for project reproducibility.
\end{enumerate}

\section*{Final Deliverables}
\begin{itemize}
    \item \textbf{Frontend:} A user-friendly website with D3.js visualizations and a form for fraud prediction.
    \item \textbf{Backend:} A REST API for fraud detection.
    \item \textbf{Deployment:} A fully functional hosted website accessible to users.
    \item \textbf{Documentation:} Detailed explanations for reproducibility and maintenance.
\end{itemize}

\end{document}
